\chapter{CONCLUSION}
\label{ch:conclusion}

\section{Summary of Contributions}

This thesis presented two complementary contributions to the design of FHE-friendly Transformer inference: the \emph{Learnable Polynomial Approximation Network (LPAN)}, which replaces all non-linear operations in BERT models with learnable polynomial approximations suitable for evaluation under CKKS homomorphic encryption, and the \emph{Synthesizer-LPAN} architectural extension, which additionally replaces the standard self-attention block with a frozen, learned, plaintext attention pattern, removing the dominant ciphertext-by-ciphertext multiplications and the polynomial Softmax entirely. The six contributions are:

\begin{enumerate}
    \item \textbf{LPAN: A Three-Stage Progressive Polynomial Replacement Framework.} A complete pipeline that replaces GELU, Softmax, and LayerNorm in BERT models through staged training: Stage~1 replaces GELU via cross-entropy fine-tuning, Stage~2 progressively replaces Softmax with per-head polynomials via attention-level knowledge distillation, and Stage~3 progressively replaces LayerNorm with co-adaptation. The framework was validated on four BERT variants (Tiny, Mini, Small, Base) on SST-2 and additionally on BERT-Base on MRPC, achieving accuracy competitive with or exceeding MPCFormer---the state-of-the-art MPC-based approach that replaces only GELU and Softmax---while replacing \emph{all three} non-linear operations non-interactively.

    \item \textbf{Distribution-Weighted Minimax Polynomial Initialisation.} A novel approximation method that incorporates the empirical activation distribution into the minimax objective, achieving $3.4\times$ lower weighted error than standard Chebyshev approximation on GELU at degree~8. These coefficients serve as informed initialisations that are subsequently refined through backpropagation.

    \item \textbf{Progressive Layer-by-Layer Replacement with Attention-Level Knowledge Distillation.} A composite loss function $\mathcal{L} = \alpha \mathcal{L}_{\mathrm{CE}} + \beta \sum \mathrm{KL}(A_{\mathrm{T}} \| A_{\mathrm{S}}) + \gamma \sum \mathrm{MSE}(H_{\mathrm{T}}, H_{\mathrm{S}})$ that preserves attention patterns and hidden representations during progressive replacement, preventing the catastrophic error amplification predicted by theoretical analysis.

    \item \textbf{Per-Head Polynomial Softmax.} An architecture where each attention head maintains independent learnable Chebyshev coefficients, enabling head specialisation for the distinct attention patterns learned by different heads.

    \item \textbf{Complete Polynomial LayerNorm Replacement.} The first approach to replace LayerNorm's inverse square root with a learnable polynomial across all layers of a deep Transformer. Co-adaptation---unfreezing all previously-replaced layers during each progressive step---proved essential for maintaining accuracy.

    \item \textbf{Synthesizer-LPAN: Architectural Elimination of $Q$, $K$, and Softmax under FHE.} An architectural extension (Section~\ref{sec:synthesizer_lpan}) that replaces the standard self-attention block with a frozen, learned, plaintext attention pattern $A_h \in \mathbb{R}^{L\times L}$ in the spirit of Tay et al.~\cite{tay2020synthesizer}. Under FHE this eliminates the $W_q$ and $W_k$ projections, the data-dependent $Q\!\cdot\!K^\top$ ciphertext-by-ciphertext multiplications, and the depth-12 polynomial Softmax \emph{entirely}. Combined with Halevi--Shoup baby-step giant-step (BSGS) fusion of the cyclic-shift mask with each plaintext diagonal~\cite{halevi2014bsgs} ($256 \to 32$ rotations per attention block), batched LPAN polynomial evaluation ($N_{\mathrm{BATCH}} = 16$), and a tuned CKKS chain ($\mathrm{chain} = 22$), Synthesizer-LPAN delivers a \emph{single-GPU sub-100-second end-to-end coherent FHE BERT result on plain HEonGPU}: \textbf{60.9\,s} for a 12-layer forward pass at $L = 128$ on a single NVIDIA H100, a $\mathbf{13.67\times}$ speed-up over the honest LPAN baseline ($833$\,s), measured on the vendored HEonGPU~\cite{kim2024heongpu} CKKS backend. Cerium~\cite{cerium2025} is faster on the same hardware via compiler/runtime optimization, so the two approaches are best viewed as complementary and in principle composable.
\end{enumerate}

\section{Key Results}

The LPAN framework achieved the following results on SST-2 and MRPC:

\begin{table}[ht]
\centering
\caption{Summary of LPAN results. BERT-Base (SST-2 and MRPC) trained on RunPod H100 SXM; all others on NVIDIA RTX~5070~Ti laptop GPU.}
\label{tab:conclusion_results}
\begin{tabular}{lcccc}
\toprule
\textbf{Model / Task} & \textbf{Baseline} & \textbf{LPAN} & \textbf{$\Delta$} & \textbf{MPCFormer} \\
\midrule
BERT-Tiny (SST-2)      & 83.26\% & 83.14\%               & $-$0.11\% & --- \\
BERT-Mini (SST-2)      & 87.16\% & 86.81\%               & $-$0.34\% & --- \\
BERT-Small (SST-2)     & 87.73\% & 88.53\%               & $+$0.80\% & --- \\
BERT-Base (SST-2)      & 93.12\% & 91.40\%               & $-$1.72\% & 91.3\% \\
BERT-Base (MRPC, F1)   & $\approx$88.9\% & 87.71\%  & $-$1.19\% & 82.1\% acc \\
\bottomrule
\end{tabular}
\end{table}

On SST-2, LPAN BERT-Base reaches 91.40\%---above MPCFormer's 91.3\%---while replacing all three non-linearities and requiring zero server--server communication. On MRPC, LPAN achieves 82.35\% accuracy / 87.71\% F1, again exceeding MPCFormer (82.1\%), confirming cross-task generality. On the FHE wall-time front, the Synthesizer-LPAN extension achieves a single-H100 end-to-end forward latency of \textbf{60.9\,s} ($N_{\mathrm{BATCH}} = 16$, $L = 128$, chain $22$), a $\mathbf{13.67\times}$ speed-up over the $833$\,s honest LPAN baseline (Section~\ref{sec:slpan_results}), and a single-GPU sub-100-second pure-FHE BERT result on plain HEonGPU. The current positioning is therefore architectural novelty plus honest baseline speedup, not current latency SOTA.

\section{Limitations}

The current work has several limitations that should be acknowledged:

\begin{enumerate}
    \item \textbf{Partial multi-task evaluation.} Full GLUE evaluation (QNLI, QQP, additional model--task pairs) is pending. The two completed tasks (SST-2, MRPC) demonstrate cross-task generality but do not yet cover natural language inference or question answering.

    \item \textbf{Encrypted inference scale.} The full CKKS encrypted forward pass for the Stage-3 LPAN BERT-Tiny was previously demonstrated as a proof-of-concept; the headline Synthesizer-LPAN $60.9$\,s number on BERT-Base is on a fresh, untrained random-weight checkpoint as a wall-time-only benchmark. The downstream accuracy of the Stage-4 Synthesizer-LPAN model on the GLUE tasks is pending and reported as \textit{TBD} in Tables~\ref{tab:slpan_glue} and~\ref{tab:slpan_pfsr_accuracy}.

    \item \textbf{Training cost.} The progressive layer-by-layer strategy for BERT-Base requires substantial GPU time for Stage~3. BERT-Base SST-2 and MRPC Stage~3 required a dedicated RunPod H100 SXM pod to complete within a reasonable time budget.

    \item \textbf{Polynomial degree sensitivity.} The depth-adaptive polynomial degree selection is currently manual. An automated approach that jointly optimises degrees across all layers and stages would improve usability.

    \item \textbf{Approximation interval sensitivity.} The profiled intervals are dataset-specific. Out-of-distribution inputs could produce activations outside the approximation interval, leading to large polynomial extrapolation errors.
\end{enumerate}

\section{Future Work}

Several directions extend this work:

\begin{enumerate}
    \item \textbf{Complete GLUE benchmark evaluation for Synthesizer-LPAN.} Extend the four-stage Synthesizer-LPAN training pipeline (Section~\ref{sec:slpan_training}) to QNLI, QQP, and additional tasks to populate the placeholder accuracy tables of Section~\ref{sec:slpan_glue}.

    \item \textbf{Composing Synthesizer-LPAN with multi-GPU FHE compilers.} The Synthesizer-LPAN circuit is, by construction, a strictly cheaper instantiation of the standard pure-FHE BERT pipeline targeted by CERIUM~\cite{cerium2025}. Re-targeting the Synthesizer-LPAN circuit through the CERIUM compiler on $8\times$B200 hardware is the most direct path to single-digit-second end-to-end FHE BERT inference.

    \item \textbf{CPU-optimised FHE inference with OpenFHE + HEXL.} Port the encrypted inference pipeline to OpenFHE with Intel HEXL AVX-512 acceleration on a dedicated 32-vCPU CPU pod to reduce single-sample latency on commodity hardware.

    \item \textbf{Larger architectures.} Apply LPAN to GPT-2, LLaMA, or other decoder-only architectures to demonstrate generality beyond encoder-only BERT.

    \item \textbf{Bootstrapping integration.} Combine LPAN's low-depth polynomial circuits with CKKS bootstrapping to enable encrypted inference on arbitrarily deep models without increasing the initial modulus size.

    \item \textbf{Automated degree and interval selection.} Develop a differentiable or search-based method to jointly optimise polynomial degrees, approximation intervals, and Chebyshev coefficients across all layers and stages.
\end{enumerate}
